\documentclass[10pt,a4paper]{article}

% Черновик для правки названий. Кинематика.

% ---------- шрифты и язык ----------
\usepackage{cmap}
\usepackage[T2A]{fontenc}
\usepackage[utf8]{inputenc}
\usepackage[russian]{babel}
\usepackage{amsmath,amssymb}
\usepackage{tempora}

% ---------- поля ----------
\usepackage[a4paper,left=14mm,right=14mm,top=14mm,bottom=14mm]{geometry}

% ---------- таблицы ----------
\usepackage{array,booktabs,longtable}
\usepackage{ragged2e}
\usepackage[table]{xcolor}

\definecolor{rule}{HTML}{C9CDC7}
\arrayrulecolor{rule}
\renewcommand{\arraystretch}{1.5}
\setlength{\parindent}{0pt}
\flushbottom

\begin{document}

{\LARGE\bfseries Кинематика}\par
\vspace{2mm}
{\large Названия карточек}\par
\vspace{3mm}
{\small Формулы перенесены из шпаргалки \texttt{09-kinematika.tex}.
Правьте второй столбец; формулы тоже можно менять.}\par
\vspace{6mm}

\begin{longtable}{@{}>{\RaggedRight\arraybackslash}p{0.46\linewidth}@{\hspace{8mm}}>{\RaggedRight\arraybackslash}p{0.46\linewidth}@{}}
\toprule
\textbf{Формула} & \textbf{Название карточки} \\
\midrule
\endfirsthead
\toprule
\textbf{Формула} & \textbf{Название карточки} \\
\midrule
\endhead
\bottomrule
\endfoot
$\displaystyle \vec v=\dfrac{\Delta\vec r}{\Delta t}$ & Мгновенная скорость \\
\addlinespace[2mm]
$\displaystyle \langle\vec v\rangle=\dfrac{\Delta\vec r}{\Delta t}$ & Средняя скорость перемещения \\
\addlinespace[2mm]
$\displaystyle \langle v\rangle=\dfrac{s}{\Delta t}$ & Средняя путевая скорость \\
\midrule
\addlinespace[2mm]
$\displaystyle \vec a=\dfrac{\Delta\vec v}{\Delta t}$ & Ускорение \\
\midrule
\addlinespace[2mm]
$\displaystyle \Delta\vec r_{\text{абс}}=\Delta\vec r_{\text{отн}}+\Delta\vec r_{\text{пер}}$ & Закон сложения перемещений \\
\addlinespace[2mm]
$\displaystyle \vec v_{\text{абс}}=\vec v_{\text{отн}}+\vec v_{\text{пер}}$ & Закон сложения скоростей \\
\addlinespace[2mm]
$\displaystyle \vec a_{\text{абс}}=\vec a_{\text{отн}}+\vec a_{\text{пер}}$ & Закон сложения ускорений \\
\midrule
\addlinespace[2mm]
$\displaystyle \vec v=\vec v_0+\vec a t$\newline
$\displaystyle \vec r=\vec r_0+\vec v_0 t+\dfrac{\vec a t^2}{2}$ & Кинематические уравнения \\
\addlinespace[2mm]
$\displaystyle \text{Ox: } a_x=\text{const}$\newline
$\displaystyle v_x=v_{0x}+a_x t$\newline
$\displaystyle x=x_0+v_{0x}t+\dfrac{a_x t^2}{2}$\newline
$\displaystyle \text{Oy: } a_y=\text{const}$\newline
$\displaystyle v_y=v_{0y}+a_y t$\newline
$\displaystyle y=y_0+v_{0y}t+\dfrac{a_y t^2}{2}$ & Кинематические уравнения (в проекциях) \\
\addlinespace[2mm]
$\displaystyle \Delta r_x=\dfrac{v_x^2-v_{0x}^2}{2a_x}$ & Перемещение (без времени) \\
\addlinespace[2mm]
$\displaystyle \Delta r_x=\dfrac{v_{0x}+v_x}{2}\,t$ & Перемещение (через среднюю скорость) \\
\midrule
\addlinespace[2mm]
$\displaystyle \text{Ox: } a_x=0$\newline
$\displaystyle v_x=v_0\cos\alpha$\newline
$\displaystyle x=x_0+v_0\cos\alpha\,t$\newline
$\displaystyle \text{Oy: } a_y=-g$\newline
$\displaystyle v_y=v_0\sin\alpha-gt$\newline
$\displaystyle y=y_0+v_0\sin\alpha\,t-\dfrac{g t^2}{2}$ & Бросок под углом к горизонту \\
\addlinespace[2mm]
$\displaystyle t_{\text{под}}=\dfrac{v_0\sin\alpha}{g}$ & Время подъема \\
\addlinespace[2mm]
$\displaystyle H=\dfrac{v_{0y}^2}{2g}$ & Высота подъема ($y_0=0$) \\
\addlinespace[2mm]
$\displaystyle t_{\text{пад}}=\sqrt{\dfrac{2H}{g}}$ & Время падения \\
\addlinespace[2mm]
$\displaystyle L=\dfrac{v_0^2\sin 2\alpha}{g}$ & Дальность полета \\
\midrule
\addlinespace[2mm]
$\displaystyle s=R\Delta\varphi$ & Длина дуги окружности \\
\addlinespace[2mm]
$\displaystyle \omega=\dfrac{\Delta\varphi}{\Delta t}$ & Угловая скорость \\
\addlinespace[2mm]
$\displaystyle \omega=\dfrac{2\pi}{T}$ & Угловая скорость (через период) \\
\addlinespace[2mm]
$\displaystyle \omega=2\pi\nu$ & Угловая скорость (через частоту) \\
\addlinespace[2mm]
$\displaystyle \nu=\dfrac{1}{T}$ & Частота и период \\
\addlinespace[2mm]
$\displaystyle v=\omega R$ & Линейная скорость (при вращательном движении) \\
\addlinespace[2mm]
$\displaystyle a=\dfrac{v^2}{R}=\omega^2 R$ & Центростремительное ускорение \\
\addlinespace[2mm]
\end{longtable}

\vspace{4mm}
{\footnotesize Всего карточек: 23.}

\end{document}
